\chapter{Modeling with Flows and Cuts}

The previous chapters built the machinery: augmenting paths, push--relabel, min-cost flow, bipartite and general matching. This chapter is about the harder, unautomatable part -- \emph{seeing} that a problem is secretly a flow problem, a cut problem, or a matching problem, and then bending it into the exact shape an off-the-shelf solver expects. The algorithms are fixed. The work is in the reduction.

A reduction is a translation. You take a problem stated in the language of projects and prerequisites, or pixels and labels, or teams and remaining games, and you produce a graph whose maximum flow, minimum cut, or maximum matching answers the original question. The translation has to hold in both directions: every legal solution of the original must correspond to a flow/cut of the graph with a matching objective value, and conversely. Getting both directions right is where reductions live or die.

\begin{focusbox}
We assume you already know how to \emph{compute} a max flow, a min cut, and a maximum matching, and that you know the max-flow min-cut theorem. This chapter never re-proves those. It proves that particular \emph{models} are correct: that the min cut of the graph we build equals the cost of the best project portfolio, the best image segmentation, the maximum set of disjoint paths, and so on. The recurring proof obligation is a two-sided correspondence between feasible solutions and flows/cuts.
\end{focusbox}

We will collect a small toolbox of gadgets -- vertex capacities, lower bounds, antiparallel edges -- then spend the bulk of the chapter on five structurally different modeling wins: project selection / maximum-weight closure, image segmentation as an $s$--$t$ min cut, baseball elimination, and edge- and vertex-disjoint paths via Menger's theorem. Each gets a hand-traced small graph.

\section{Conventions and the modeling reflex}
\label{ch18:conventions}

Fix notation for the whole chapter. A flow network is a directed graph $\Gph=(V,E)$ with a \emph{source} $s$, a \emph{sink} $t$, and a nonnegative integer (or rational) capacity $c(u,v)$ on each edge. A flow $f$ assigns $0 \le f(u,v) \le c(u,v)$ to each edge and conserves flow at every vertex except $s$ and $t$. Its value $\abs{f}$ is the net flow out of $s$. An $s$--$t$ cut $(S,T)$ partitions $V$ with $s\in S$, $t\in T$; its capacity is $c(S,T)=\sum_{u\in S, v\in T} c(u,v)$, counting only edges crossing \emph{forward} from $S$ to $T$. Max-flow min-cut says $\max_f \abs{f} = \min_{(S,T)} c(S,T)$.

\begin{intuitionbox}
Three reflexes cover almost every modeling problem.
\begin{itemize}
\item \textbf{``Choose a subset, pay for crossings''} $\Rightarrow$ \emph{minimum cut}. A cut \emph{is} a partition of the vertices into two sides. If your problem is ``put each thing on one of two sides and minimize the cost of pairs that end up split,'' it is a cut problem in disguise.
\item \textbf{``Route as much / as many as possible through a bottlenecked system''} $\Rightarrow$ \emph{maximum flow}. Capacities are the bottlenecks; the value is the throughput.
\item \textbf{``Pair up / cover / select with conflicts''} $\Rightarrow$ \emph{matching} (a special flow). Disjointness constraints become unit capacities.
\end{itemize}
The whole game is choosing the vertices, the edges, and especially the capacities so that the solver's optimum is forced to coincide with yours.
\end{intuitionbox}

\subsection{Gadget 1: vertex capacities}
\label{ch18:vertex-cap}

Standard flow puts capacities on edges. Many problems want to limit how much passes \emph{through a vertex} -- a router that can forward at most $k$ packets, an intersection that can carry at most $c$ cars. The fix is vertex splitting.

\begin{focusbox}
\textbf{Vertex split.} Replace each capacitated vertex $v$ by two copies $v_{\text{in}}$ and $v_{\text{out}}$ joined by a single edge $v_{\text{in}}\to v_{\text{out}}$ with capacity equal to the vertex capacity $c(v)$. Every original edge $(u,v)$ becomes $(u_{\text{out}}, v_{\text{in}})$; every $(v,w)$ becomes $(v_{\text{out}}, w_{\text{in}})$.
\end{focusbox}

\begin{center}
\begin{tikzpicture}[>=Stealth, node distance=14mm, every node/.style={font=\small}]
  \node[circle,draw,fill=nodeblue!15] (u) {$u$};
  \node[circle,draw,fill=nodeblue!15,right=18mm of u] (v) {$v$};
  \node[circle,draw,fill=nodeblue!15,right=18mm of v] (w) {$w$};
  \draw[->] (u) -- (v);
  \draw[->] (v) -- (w);
  \node[below=2mm of v] {capacity $c(v)$ on the \emph{vertex}};

  \begin{scope}[yshift=-26mm]
    \node[circle,draw,fill=nodeblue!15] (u2) {$u_{\text{o}}$};
    \node[circle,draw,fill=quotientorange!25,right=14mm of u2] (vi) {$v_{\text{i}}$};
    \node[circle,draw,fill=quotientorange!25,right=12mm of vi] (vo) {$v_{\text{o}}$};
    \node[circle,draw,fill=nodeblue!15,right=14mm of vo] (w2) {$w_{\text{i}}$};
    \draw[->] (u2) -- node[above,font=\scriptsize]{$\infty$} (vi);
    \draw[->,very thick,quotientorange] (vi) -- node[above,font=\scriptsize]{$c(v)$} (vo);
    \draw[->] (vo) -- node[above,font=\scriptsize]{$\infty$} (w2);
  \end{scope}
\end{tikzpicture}
\end{center}

\begin{whybox}
Every unit of flow that visits $v$ in the original must traverse the unique internal edge $v_{\text{in}}\to v_{\text{out}}$ in the split graph, because that edge is the \emph{only} way to get from any $v_{\text{in}}$ to any $v_{\text{out}}$. So the through-flow at $v$ is exactly the flow on that edge, which is capped at $c(v)$. Conversely, any flow in the split graph projects to a flow in the original obeying the vertex bound. The two networks have identical max-flow values, and a min cut that uses an internal edge corresponds to deleting the vertex. The construction at most doubles $\abs{V}$ and adds $\abs{V}$ edges, so it is free asymptotically.
\end{whybox}

\subsection{Gadget 2: lower bounds on edges}
\label{ch18:lower-bounds}

Sometimes an edge must carry \emph{at least} $\ell(u,v)$ units, not just at most $c(u,v)$. Think of a delivery contract that obligates a minimum shipment, or a scheduling constraint that a machine run at least some minimum load. A flow with lower bounds need not even exist, so the first question is feasibility.

\begin{focusbox}
\textbf{Demand reduction.} Given bounds $\ell(u,v) \le f(u,v) \le c(u,v)$, substitute $f = \ell + g$ where $0 \le g(u,v) \le c(u,v)-\ell(u,v)$. Conservation of $f$ becomes conservation of $g$ plus a fixed \emph{excess} at each vertex:
\[
  d(v) \;=\; \sum_{u} \ell(u,v) \;-\; \sum_{w} \ell(v,w),
\]
the surplus of forced inflow over forced outflow. We must source this surplus or sink this deficit.
\end{focusbox}

\begin{algorithm}[H]
\caption{Feasible circulation with lower bounds}
\label{ch18:alg-lowerbound}
\DontPrintSemicolon
\KwIn{Digraph with bounds $\ell(u,v)\le f(u,v)\le c(u,v)$.}
\KwOut{A feasible flow $f$, or report infeasible.}
Build new graph on the same vertices; add super-source $S$ and super-sink $T$.\;
\For{each edge $(u,v)$}{
  Add edge $(u,v)$ with capacity $c(u,v)-\ell(u,v)$.\;
}
\For{each vertex $v$}{
  Compute $d(v) = \sum_u \ell(u,v) - \sum_w \ell(v,w)$.\;
  \lIf{$d(v) > 0$}{add $S\to v$ with capacity $d(v)$}
  \lIf{$d(v) < 0$}{add $v\to T$ with capacity $-d(v)$}
}
Compute max flow $S\rightsquigarrow T$.\;
\eIf{every $S$-edge is saturated (flow $=\sum_{d(v)>0} d(v)$)}{
  \Return $f(u,v) = \ell(u,v) + g(u,v)$ for each original edge.\;
}{
  \Return infeasible.\;
}
\end{algorithm}

\begin{whybox}
The super-source $S$ injects exactly the forced surplus into each deficit-of-outflow vertex, and the super-sink $T$ absorbs exactly the forced excess elsewhere. A feasible $g$ exists iff $S$ can push all of its supply to $T$, i.e.\ iff the max flow saturates every $S$-edge. Saturation is equivalent to value $=\sum_{d(v)>0}d(v)$. Reconstructing $f=\ell+g$ restores the lower bounds while preserving the upper bounds because $g \le c-\ell$. To turn this from a feasibility test into a max-flow-with-lower-bounds (between a real $s$ and $t$), add an edge $t\to s$ of capacity $\infty$ to convert the problem to a circulation, find a feasible circulation, then maximize the residual $s\rightsquigarrow t$ flow on top.
\end{whybox}

\begin{center}
\begin{tikzpicture}[>=Stealth, every node/.style={font=\small}, node distance=14mm]
  \node[circle,draw,fill=nodeblue!15] (n1) {$1$};
  \node[circle,draw,fill=nodeblue!15, right=20mm of n1] (n2) {$2$};
  \node[circle,draw,fill=nodeblue!15, below=20mm of n2] (n3) {$3$};
  \node[circle,draw,fill=nodeblue!15, below=20mm of n1] (n4) {$4$};
  \draw[->] (n1) -- node[above,font=\scriptsize]{$(2,4)$} (n2);
  \draw[->] (n2) -- node[right,font=\scriptsize]{$(1,3)$} (n3);
  \draw[->] (n3) -- node[below,font=\scriptsize]{$(2,2)$} (n4);
  \draw[->] (n4) -- node[left,font=\scriptsize]{$(0,5)$} (n1);
  \node[circle,draw,fill=sourceblue!25, left=14mm of n1, yshift=-10mm] (S) {$S$};
  \node[circle,draw,fill=targetgreen!25, right=14mm of n2, yshift=-10mm] (T) {$T$};
  \draw[->,sourceblue] (S) to[bend left=10] node[above,font=\scriptsize]{$1$} (n2);
  \draw[->,sourceblue] (S) -- node[below,font=\scriptsize]{$2$} (n4);
  \draw[->,targetgreen] (n1) -- node[above,font=\scriptsize]{$2$} (T);
  \draw[->,targetgreen] (n3) to[bend right=10] node[below,font=\scriptsize]{$1$} (T);
\end{tikzpicture}
\end{center}
\vspace{-0.4em}
\begin{labbox}
A four-cycle circulation, no real source or sink (the diagram above shows the cycle with $(\ell,c)$ labels and the super-source/sink attachments computed below). Vertices $\{1,2,3,4\}$; edges with $(\ell,c)$: $1\to2$ is $(2,4)$, $2\to3$ is $(1,3)$, $3\to4$ is $(2,2)$, $4\to1$ is $(0,5)$. Every vertex has one incoming and one outgoing edge, so compute each $d(v)=(\text{forced in})-(\text{forced out})$:
\[
  d(1) = \ell(4{\to}1)-\ell(1{\to}2) = 0-2 = -2,\quad
  d(2) = 2-1 = +1,
\]
\[
  d(3) = 1-2 = -1,\quad
  d(4) = 2-0 = +2.
\]
The positive-$d$ vertices ($2$ and $4$) get edges \emph{from} the super-source $S$ with capacities $1$ and $2$; the negative-$d$ vertices ($1$ and $3$) get edges \emph{to} the super-sink $T$ with capacities $2$ and $1$. Total supply $\sum_{d>0} d = 1+2 = 3$, total demand $\sum_{d<0}|d| = 2+1 = 3$ -- they always match, which is a useful arithmetic sanity check. The residual capacities $c-\ell$ on the original edges are $1\to2:2$, $2\to3:2$, $3\to4:0$, $4\to1:5$.
\par
Now push $S$'s $3$ units to $T$. The wiring is $S\to2$ (cap $1$), $S\to4$ (cap $2$) on the supply side, and $1\to T$ (cap $2$), $3\to T$ (cap $1$) on the demand side, with the four original edges carrying their residual capacities. The edge $3\to4$ has residual capacity $0$, so the route through it is closed, but the remaining residual edges suffice. A feasible routing: $S\to4\to1\to T$ carries $2$ units (uses residual $4\to1=5$, delivers $2$ to vertex $1$'s sink edge, exactly its cap), and $S\to2\to3\to T$ carries $1$ unit (uses residual $2\to3=2$, delivers to vertex $3$'s cap-$1$ sink edge). All $3$ source units placed, every $S$-edge saturated: \textbf{feasible}. Reconstruct $f=\ell+g$: $f(1{\to}2)=2+0=2$, $f(2{\to}3)=1+1=2$, $f(3{\to}4)=2+0=2$, $f(4{\to}1)=0+2=2$. Check conservation: every vertex carries $2$ in and $2$ out. A uniform circulation of value $2$ satisfies all the lower bounds.
\end{labbox}

\subsection{Gadget 3: antiparallel and undirected edges}
\label{ch18:antiparallel}

Two edges $(u,v)$ and $(v,u)$ both present (``antiparallel'') confuse the textbook residual-graph bookkeeping, which wants a single residual capacity per ordered pair. Split one of them: replace $(u,v)$ by $u\to x\to v$ through a fresh vertex $x$, each new edge inheriting $c(u,v)$. For an \emph{undirected} edge of capacity $c$, model it as two directed edges $(u,v)$ and $(v,u)$ each of capacity $c$; a flow uses at most one direction at full strength in any optimal solution, and the cut capacity counts the edge once per crossing direction -- exactly right, since an undirected edge crossing a cut should be paid for once.

\begin{trapbox}
For \emph{undirected min cut}, both directed copies have capacity $c$, and a cut that separates the endpoints pays $c$ (one forward copy crosses). Do \emph{not} also count the backward copy: in an $s$--$t$ cut $(S,T)$ only $S\to T$ edges are charged, so the $v\to u$ copy with $v\in T, u\in S$ is a backward edge and contributes $0$. The single undirected edge is therefore charged exactly $c$. People who symmetrize sloppily double every cut.
\end{trapbox}

\section{Maximum-weight closure and project selection}
\label{ch18:closure}

Now the first substantial reduction. You have a set of projects. Each project $i$ has a profit $p_i$ (possibly negative -- it might be a pure cost). Some projects \emph{require} others: to do $i$ you must also do every project it depends on. Choose a subset maximizing total profit subject to all prerequisites being satisfied. This is the canonical ``open-pit mining'' / ``task selection'' problem, and it is a minimum cut.

\begin{definition}[Closure]
\label{ch18:def-closure}
Let $\Gph=(V,E)$ be a directed graph with vertex weights $w:V\to\R$. A subset $C\subseteq V$ is \emph{closed} if every edge leaving $C$ stays inside $C$: whenever $u\in C$ and $(u,v)\in E$, then $v\in C$. The \emph{maximum-weight closure} maximizes $\sum_{v\in C} w(v)$ over closed $C$.
\end{definition}

Read the edges as ``$u$ needs $v$.'' Closure says: if you take $u$, you take everything $u$ points to. That is exactly the prerequisite rule.

\begin{intuitionbox}
Picture two heaps: projects you \emph{select} (side $S$, with the source) and projects you \emph{reject} (side $T$, with the sink). Profitable projects ``want'' to be on the source side; costly ones want the sink side. A prerequisite edge $u\to v$ must never go from $S$ to $T$ (that would mean taking $u$ but not its prerequisite $v$). Give that edge infinite capacity so no finite cut dares to sever it. Then minimizing the cut balances the forgone profit of rejected good projects against the incurred cost of accepted bad projects -- which is precisely the negation of total profit.
\end{intuitionbox}

\subsection{The construction}

\begin{focusbox}
Build a flow network. Add source $s$ and sink $t$.
\begin{itemize}
\item For each vertex $v$ with $w(v) > 0$: edge $s\to v$ with capacity $w(v)$.
\item For each vertex $v$ with $w(v) < 0$: edge $v\to t$ with capacity $-w(v)$.
\item For each prerequisite edge $(u,v)\in E$: edge $u\to v$ with capacity $\infty$.
\end{itemize}
Let $P = \sum_{v:\,w(v)>0} w(v)$ be the total positive weight. Then the maximum closure value equals $P - (\text{min } s\text{--}t\text{ cut})$, and the selected set is $S\setminus\{s\}$ for a min cut $(S,T)$.
\end{focusbox}

\begin{center}
\begin{tikzpicture}[>=Stealth, every node/.style={font=\small}, node distance=12mm]
  \node[circle,draw,fill=sourceblue!25] (s) {$s$};
  \node[circle,draw,fill=nodeblue!15, above right=4mm and 22mm of s] (a) {$a$};
  \node[circle,draw,fill=nodeblue!15, right=22mm of s] (b) {$b$};
  \node[circle,draw,fill=nodeblue!15, below right=4mm and 22mm of s] (c) {$c$};
  \node[circle,draw,fill=targetgreen!25, right=22mm of b] (t) {$t$};
  \draw[->] (s) -- node[above,font=\scriptsize]{$3$} (a);
  \draw[->] (s) -- node[above,font=\scriptsize]{$5$} (b);
  \draw[->] (c) -- node[below,font=\scriptsize]{$4$} (t);
  \draw[->,structurepurple] (a) to[bend left=10] node[above,font=\scriptsize]{$\infty$} (c);
  \draw[->,structurepurple] (b) -- node[above,font=\scriptsize,pos=.4]{$\infty$} (c);
  \node[below=10mm of b,font=\scriptsize,align=center] {$w(a)=3,\ w(b)=5,\ w(c)=-4$.\\ Edges $a\to c$, $b\to c$ mean $a$ and $b$ each need $c$.};
\end{tikzpicture}
\end{center}

\begin{theorem}[Closure as min cut]
\label{ch18:thm-closure}
With the construction above, an $s$--$t$ cut has \emph{finite} capacity iff its source side $S\setminus\{s\}$ is a closed set, and in that case
\[
  c(S,T) \;=\; P \;-\; \sum_{v\in S\setminus\{s\}} w(v).
\]
Consequently the minimum cut yields the maximum-weight closure, with value $P - c(S,T)_{\min}$.
\end{theorem}

\begin{proof}
Let $C = S\setminus\{s\}$. A cut is finite iff no infinite edge crosses forward. The infinite edges are the prerequisite edges $u\to v$; one crosses forward iff $u\in C$ and $v\notin C$. So finiteness is exactly the closure condition: every prerequisite of a selected vertex is selected.

Now compute the finite cut capacity. The only finite edges are the $s\to v$ edges (capacity $w(v)$ for $w(v)>0$) and the $v\to t$ edges (capacity $-w(v)$ for $w(v)<0$).
\begin{itemize}
\item An edge $s\to v$ crosses forward iff $v\in T$, i.e.\ $v\notin C$. It contributes $w(v)$. Summing over positive-weight $v\notin C$ gives $\sum_{v\notin C,\, w(v)>0} w(v) = P - \sum_{v\in C,\, w(v)>0} w(v)$.
\item An edge $v\to t$ crosses forward iff $v\in S$, i.e.\ $v\in C$. It contributes $-w(v)$. Summing over negative-weight $v\in C$ gives $\sum_{v\in C,\, w(v)<0} (-w(v)) = -\sum_{v\in C,\,w(v)<0} w(v)$.
\end{itemize}
Adding,
\[
  c(S,T) = P - \sum_{v\in C,\, w(v)>0} w(v) - \sum_{v\in C,\, w(v)<0} w(v) = P - \sum_{v\in C} w(v).
\]
Since $P$ is a constant, minimizing $c(S,T)$ maximizes $\sum_{v\in C} w(v)$ over closed $C$. The infinite edges guarantee the minimizer is closed.
\end{proof}

\begin{labbox}
Run it on the picture. $P = w(a)+w(b) = 3+5 = 8$. The min $s$--$t$ cut: the only finite edges are $s\to a$ (3), $s\to b$ (5), $c\to t$ (4). Options for a closed set $C$:
\begin{itemize}
\item $C=\emptyset$: cut cuts $s\to a, s\to b$, capacity $8$. Closure value $8-8=0$.
\item $C=\{a,b,c\}$: closed (both prerequisites of $c$ present). Cut cuts only $c\to t$, capacity $4$. Closure value $8-4=4$.
\item $C=\{a,c\}$: closed. Cut cuts $s\to b$ (5) and $c\to t$ (4), capacity $9$. Value $8-9=-1$.
\end{itemize}
Minimum cut is $4$, achieved by $C=\{a,b,c\}$: take all three, profit $3+5-4 = 4$. The cut $4$ matches $P - 4 = 4$. Note $C=\{a\}$ is \emph{not} closed ($a\to c$ leaves), so the infinite edge $a\to c$ forbids it -- any cut separating $a$ from $c$ is infinite.
\end{labbox}

\begin{algorithm}[H]
\caption{Project selection via maximum-weight closure}
\label{ch18:alg-projsel}
\DontPrintSemicolon
\KwIn{Projects $1..n$ with profits $p_i$; prerequisite pairs $(i,j)$ meaning $i$ needs $j$.}
\KwOut{Subset maximizing total profit subject to prerequisites.}
$P \gets 0$; create source $s$, sink $t$.\;
\For{$i=1$ \KwTo $n$}{
  \lIf{$p_i>0$}{add $s\to i$ with capacity $p_i$; \ $P \mathrel{+}= p_i$}
  \lElseIf{$p_i<0$}{add $i\to t$ with capacity $-p_i$}
}
\lForEach{prerequisite $(i,j)$}{add $i\to j$ with capacity $\infty$}
$(F,\, (S,T)) \gets \textsc{MaxFlowMinCut}(s,t)$.\;
\Return $S\setminus\{s\}$ as the selected projects; profit $= P - \abs{F}$.\;
\end{algorithm}

\begin{complexitybox}
The network has $n+2$ vertices and at most $n + \abs{E}$ edges. With Dinic's algorithm the cost is $\Oh\!\big((n+\abs{E})\cdot \abs{V}^2\big)$ in general, but on unit-ish or bounded-capacity instances it is far better in practice. Memory is $\Oh(n+\abs{E})$. The $\infty$ capacities can be taken as $1 + \sum_i \abs{p_i}$ to keep arithmetic bounded -- no finite optimal cut ever uses one.
\end{complexitybox}

\begin{historybox}
The maximum-closure model is due to Jean-Claude Picard (1976), building on the open-pit mining formulation of Lerchs and Grossmann (1965). Mining engineers wanted to decide which blocks of ore to excavate: digging a deep block forces removing the blocks above it (the prerequisites), each block has a value (ore revenue minus extraction cost), and the optimal pit is a maximum-weight closure. The same template reappears as ``task scheduling with precedence,'' ``flow-based feature selection,'' and Markov-random-field energy minimization.
\end{historybox}

\begin{applicationbox}
Beyond mining: choosing which experiments to fund when some share expensive shared equipment; selecting a profitable set of software features where features depend on shared libraries; and the densest-subgraph and maximum-density problems, which reduce to a parametric family of closure/min-cut instances (Exercise~\ref{ch18:ex-densest}).
\end{applicationbox}

\section{Image segmentation as an \texorpdfstring{$s$--$t$}{s-t} min cut}
\label{ch18:segmentation}

Foreground/background segmentation: each pixel must be labeled \emph{object} or \emph{background}. Two forces compete. A \emph{data} term says how much each pixel, on its own evidence (color, intensity), prefers each label. A \emph{smoothness} term penalizes neighboring pixels that get different labels, because real objects are mostly contiguous. Minimizing data $+$ smoothness is, for two labels, exactly a minimum cut.

\begin{intuitionbox}
Put the source $s$ on the ``object'' side and the sink $t$ on the ``background'' side. A cut assigns every pixel to one side -- a labeling. Edge $s\to p$ with weight $a_p$ = ``cost of \emph{not} calling $p$ object'' (paid when $p$ lands on the background side). Edge $p\to t$ with weight $b_p$ = ``cost of \emph{not} calling $p$ background.'' A neighbor edge $p$--$q$ with weight $\lambda_{pq}$ is paid exactly when $p$ and $q$ are split. The min cut minimizes total penalty: the optimal labeling.
\end{intuitionbox}

\subsection{Construction and energy}

Let the pixels be vertices. Choose data costs: $a_p \ge 0$ is the penalty for labeling $p$ background, $b_p \ge 0$ for labeling $p$ object. For each neighboring pair $\{p,q\}$ choose a separation penalty $\lambda_{pq}\ge 0$.

\begin{focusbox}
\textbf{Segmentation network.}
\begin{itemize}
\item Source $s$ (object), sink $t$ (background).
\item Edge $s\to p$ with capacity $a_p$ for every pixel $p$.
\item Edge $p\to t$ with capacity $b_p$ for every pixel $p$.
\item For each neighbor pair $\{p,q\}$, \emph{two} directed edges $p\to q$ and $q\to p$, each capacity $\lambda_{pq}$.
\end{itemize}
A cut $(S,T)$ labels $p$ \emph{object} if $p\in S$, \emph{background} if $p\in T$.
\end{focusbox}

\begin{center}
\begin{tikzpicture}[>=Stealth, every node/.style={font=\small}]
  \node[circle,draw,fill=sourceblue!25] (s) {$s$};
  \node[circle,draw,fill=nodeblue!15, right=20mm of s, yshift=10mm] (p) {$p$};
  \node[circle,draw,fill=nodeblue!15, right=20mm of s, yshift=-10mm] (q) {$q$};
  \node[circle,draw,fill=targetgreen!25, right=20mm of p, yshift=-10mm] (t) {$t$};
  \draw[->] (s) -- node[above,font=\scriptsize]{$a_p$} (p);
  \draw[->] (s) -- node[below,font=\scriptsize]{$a_q$} (q);
  \draw[->] (p) -- node[above,font=\scriptsize]{$b_p$} (t);
  \draw[->] (q) -- node[below,font=\scriptsize]{$b_q$} (t);
  \draw[<->,quotientorange,thick] (p) -- node[right,font=\scriptsize]{$\lambda_{pq}$} (q);
\end{tikzpicture}
\end{center}

\begin{theorem}[Segmentation energy equals cut capacity]
\label{ch18:thm-seg}
For any labeling with object set $O$ (the pixels in $S$) and background set $B$ (the pixels in $T$), the cut capacity equals the energy
\[
  E(O,B) \;=\; \sum_{p\in B} a_p \;+\; \sum_{p\in O} b_p \;+\; \sum_{\{p,q\}:\, \text{split}} \lambda_{pq}.
\]
Hence the minimum cut is the minimum-energy two-label segmentation.
\end{theorem}

\begin{proof}
Tally forward-crossing edges of the cut $(S,T)$, $s\in S$, $t\in T$, $O=S\setminus\{s\}$, $B=T\setminus\{t\}$.
\begin{itemize}
\item $s\to p$ crosses iff $p\in T$, i.e.\ $p\in B$: contributes $a_p$. Total $\sum_{p\in B} a_p$.
\item $p\to t$ crosses iff $p\in S$, i.e.\ $p\in O$: contributes $b_p$. Total $\sum_{p\in O} b_p$.
\item For a neighbor pair $\{p,q\}$: if both on the same side, neither directed copy crosses forward. If split, exactly one copy crosses forward (the one going from the $S$-side endpoint to the $T$-side endpoint), contributing $\lambda_{pq}$ once.
\end{itemize}
Summing gives $E(O,B)$. Minimizing cut capacity minimizes energy.
\end{proof}

\begin{whybox}
Why does the two-edge gadget pay $\lambda_{pq}$ \emph{once} per split and \emph{zero} when not split? Because a cut only charges $S\to T$ edges. If $p,q\in S$, both copies are inside $S$, charged $0$. If $p,q\in T$, both inside $T$, charged $0$. If $p\in S, q\in T$, the copy $p\to q$ is forward (charged $\lambda_{pq}$) while $q\to p$ is backward (charged $0$). Symmetric in the other split case. This is the undirected-edge trick of Section~\ref{ch18:antiparallel} doing exactly the right accounting.
\end{whybox}

\begin{labbox}
Three pixels in a row, $p_1 - p_2 - p_3$, neighbors $\{p_1,p_2\}$ and $\{p_2,p_3\}$, all $\lambda = 2$. Data: $a = (5, 1, 0)$ (object-affinity), $b = (0, 4, 6)$ (background-affinity). So $p_1$ strongly wants object ($a_1=5$ penalty to put it in background), $p_3$ strongly wants background.
\par
Try labeling $O=\{p_1\}$, $B=\{p_2,p_3\}$. Split pairs: $\{p_1,p_2\}$ only. Energy $= (a_2+a_3) + (b_1) + \lambda = (1+0) + 0 + 2 = 3$.
\par
Try $O=\{p_1,p_2\}$, $B=\{p_3\}$. Split pairs: $\{p_2,p_3\}$. Energy $= a_3 + (b_1+b_2) + \lambda = 0 + (0+4) + 2 = 6$.
\par
Try all object: energy $= 0 + (b_1+b_2+b_3) + 0 = 10$. Try all background: $(a_1+a_2+a_3) + 0 + 0 = 6$.
\par
The winner is $O=\{p_1\}, B=\{p_2,p_3\}$ at energy $3$: $p_1$ is foreground, the rest background, with one boundary paid. A min-cut solver finds this as the cut $\{s,p_1\}\,|\,\{p_2,p_3,t\}$ of capacity $3$.
\end{labbox}

\begin{trapbox}
Min cut handles \emph{exactly two} labels. With three or more labels the energy minimization is NP-hard in general (it is multiway cut). The $\alpha$-expansion and $\alpha\beta$-swap algorithms of Boykov--Veksler--Zabih attack the multilabel case by solving a \emph{sequence} of two-label min cuts, each one a binary ``keep your label vs.\ switch to $\alpha$'' decision. They give a $2$-approximation under a metric smoothness condition, not the exact optimum (Exercise~\ref{ch18:ex-multiway}).
\end{trapbox}

\begin{trapbox}
The exact-cut reduction needs the pairwise term to be \emph{submodular}: the penalty for ``different'' must be at least the penalty for ``same.'' Here ``same'' costs $0$ and ``different'' costs $\lambda_{pq}\ge 0$, so submodularity holds automatically. If you tried to \emph{reward} disagreement (negative $\lambda$), the problem stops being a min cut and becomes NP-hard. Submodularity is the dividing line between ``min-cut easy'' and ``hard'' (Exercise~\ref{ch18:ex-seg-hard}).
\end{trapbox}

\begin{historybox}
Greig, Porteous, and Seheult (1989) first observed that the maximum a-posteriori estimate of a binary image under an Ising-type prior is a min cut -- a result that sat almost unused for a decade until Boykov and Kolmogorov (2001--2004) made min-cut/max-flow fast enough for vision and built the $\alpha$-expansion framework. Today graph cuts are a standard building block in stereo, segmentation, and image stitching.
\end{historybox}

\begin{connectionbox}
The energy $\sum_p (\text{data}) + \sum_{p\sim q}\lambda_{pq}[x_p\ne x_q]$ is, up to constants, a weighted graph \emph{cut function}, and minimizing $[x_p\ne x_q]$ summed over edges is the quadratic form $\tfrac14 x^\top \Lap x$ on $\{\pm1\}$ labels, where $\Lap$ is the graph Laplacian (Chapter~31). The data terms shift it to a linear-plus-quadratic energy. The $s$--$t$ cut solves this Boolean quadratic exactly because the pairwise part is submodular; the spectral relaxation (Chapter~32) drops the integrality and minimizes the Rayleigh quotient instead, giving Cheeger-type cuts. Same Laplacian, two different relaxations.
\end{connectionbox}

\section{Baseball elimination}
\label{ch18:baseball}

A reduction that does not look like flow at all from its statement. It is late in the season. Team $x$ has won $w_x$ games and has $r_x$ games left. Between every pair $\{i,j\}$ a known number $g_{ij}$ of games remain to be played. Question: can team $x$ \emph{still} finish in first place (tie allowed), assuming the remaining games can come out any way we like? If not, $x$ is mathematically eliminated.

\begin{intuitionbox}
Team $x$ does best if it wins all $r_x$ remaining games, reaching $W = w_x + r_x$. For $x$ to finish first, \emph{no other team} may exceed $W$ wins. So every other team $i$ can win at most $W - w_i$ more games -- call this its \emph{slack}. The remaining games among the other teams must be distributed without busting any slack. ``Distribute games subject to capacity'' is a flow: route each remaining game to one of its two teams, respecting each team's slack as a capacity into the sink.
\end{intuitionbox}

\subsection{Construction}

Fix the team $x$ we are testing. Let $W = w_x + r_x$. If some other team already has $w_i > W$, then $x$ is trivially eliminated. Otherwise build:

\begin{focusbox}
\textbf{Elimination network for team $x$.}
\begin{itemize}
\item Source $s$, sink $t$.
\item A \emph{game node} $g_{ij}$ for each unordered pair $i\ne j$ of teams other than $x$ with $g_{ij}>0$. Edge $s\to g_{ij}$ with capacity $g_{ij}$ (these games must be played and assigned).
\item A \emph{team node} for each team $i\ne x$. Edges $g_{ij}\to i$ and $g_{ij}\to j$, each capacity $\infty$ (a game between $i$ and $j$ is won by one of them).
\item Edge $i\to t$ with capacity $W - w_i$ (team $i$'s slack).
\end{itemize}
\end{focusbox}

\begin{center}
\begin{tikzpicture}[>=Stealth, every node/.style={font=\small}, node distance=8mm]
  \node[circle,draw,fill=sourceblue!25] (s) {$s$};
  \node[rectangle,draw,fill=quotientorange!20, right=16mm of s, yshift=14mm] (g12) {$g_{12}$};
  \node[rectangle,draw,fill=quotientorange!20, right=16mm of s] (g13) {$g_{13}$};
  \node[rectangle,draw,fill=quotientorange!20, right=16mm of s, yshift=-14mm] (g23) {$g_{23}$};
  \node[circle,draw,fill=nodeblue!15, right=20mm of g12, yshift=-3mm] (n1) {$1$};
  \node[circle,draw,fill=nodeblue!15, right=20mm of g13] (n2) {$2$};
  \node[circle,draw,fill=nodeblue!15, right=20mm of g23, yshift=3mm] (n3) {$3$};
  \node[circle,draw,fill=targetgreen!25, right=18mm of n2] (t) {$t$};
  \draw[->] (s) -- node[above,font=\scriptsize]{$g_{12}$} (g12);
  \draw[->] (s) -- node[above,font=\scriptsize]{$g_{13}$} (g13);
  \draw[->] (s) -- node[below,font=\scriptsize]{$g_{23}$} (g23);
  \draw[->] (g12) -- (n1); \draw[->] (g12) -- (n2);
  \draw[->] (g13) -- (n1); \draw[->] (g13) -- (n3);
  \draw[->] (g23) -- (n2); \draw[->] (g23) -- (n3);
  \draw[->] (n1) -- node[above,font=\scriptsize]{$W{-}w_1$} (t);
  \draw[->] (n2) -- node[above,font=\scriptsize]{$W{-}w_2$} (t);
  \draw[->] (n3) -- node[below,font=\scriptsize]{$W{-}w_3$} (t);
\end{tikzpicture}
\end{center}

\begin{theorem}[Baseball elimination]
\label{ch18:thm-baseball}
Let $R = \sum_{i<j,\, i,j\ne x} g_{ij}$ be the total number of remaining games among teams other than $x$. Team $x$ is \emph{not} eliminated iff the max flow in the network saturates the source, i.e.\ has value $R$ (equivalently, the min cut has capacity $R$).
\end{theorem}

\begin{proof}
($\Leftarrow$) Suppose max flow $=R$. Every game edge $s\to g_{ij}$ is saturated, so all $g_{ij}$ units leave $g_{ij}$ along $g_{ij}\to i$ and $g_{ij}\to j$. Interpret the flow on $g_{ij}\to i$ as ``$i$ wins that many of the $g_{ij}$ games.'' By conservation, the wins assigned to each team total $\sum_j f(g_{ij}\to i)$, which is the flow into team node $i$, hence at most its capacity $W-w_i$. So $i$'s final win total $w_i + (\text{assigned wins}) \le w_i + (W - w_i) = W$. Let $x$ win all its remaining games to reach $W$; then no team exceeds $W$, and $x$ ties or leads. Not eliminated.

($\Rightarrow$) Suppose $x$ is not eliminated, so there is an outcome of \emph{all} remaining games with $x$ on top at some total $\le W$ (and $x$ can always be assumed to win all its own games, pushing it to exactly $W$ and only helping). In that outcome every other team finishes with $\le W$ wins, so the games it won among the others-only games number at most $W - w_i$. Route each others-only game to its winner: this is a feasible flow saturating every $s\to g_{ij}$ edge (each game is won by someone), with team-inflows $\le W - w_i$ respecting the sink capacities. Its value is $R$. So max flow $=R$.
\end{proof}

\begin{labbox}
Three other teams $\{1,2,3\}$, current wins $w=(83, 80, 78)$. We test team $x$ with $w_x=81$, $r_x=2$, so $W = 83$. Slacks: $W-w_1 = 0$, $W-w_2 = 3$, $W-w_3 = 5$. Remaining others-only games: $g_{12}=1, g_{13}=0, g_{23}=4$, so $R = 5$.
\par
Source edges: $s\to g_{12}$ (cap 1), $s\to g_{23}$ (cap 4). Team 1 has slack $0$: edge $1\to t$ has capacity $0$. The game $g_{12}$ must send its $1$ unit to team $1$ or team $2$; team $1$ is full (capacity $0$), so the unit goes to team $2$. The game $g_{23}$ sends $4$ units split between teams $2$ and $3$. Team $2$ already took $1$ unit, leaving slack $3-1 = 2$; team $3$ has slack $5$. Together they absorb $2 + 5 = 7 \ge 4$. So all $5$ units route through: max flow $=5=R$. Team $x$ is \emph{not} eliminated.
\par
Now suppose instead $w_x = 79, r_x = 2$, so $W = 81$, slacks $(W-w_i) = (-2, 1, 3)$. Team 1 already has $83 > 81$: trivial elimination, no flow needed.
\end{labbox}

\begin{whybox}
The min cut tells you \emph{why} a team is eliminated. If the max flow is $< R$, the min cut $(S,T)$ exposes a set of teams $A = \{i : i\in S\}$ whose combined remaining games among themselves cannot fit under their combined slack. Concretely, eliminate $x$ iff there is a subset $A$ of the others with
\[
  \frac{w(A) + g(A)}{\abs{A}} \;>\; W,
\]
where $w(A)=\sum_{i\in A} w_i$ and $g(A)$ counts games with both teams in $A$: their games \emph{must} be won by members of $A$, so the average final win count of $A$ is forced above $W$. The flow finds this certificate of elimination automatically -- it is the team subset on the source side of the min cut.
\end{whybox}

\begin{algorithm}[H]
\caption{Is team $x$ eliminated?}
\label{ch18:alg-baseball}
\DontPrintSemicolon
\KwIn{Wins $w_i$, remaining pairwise games $g_{ij}$, the team $x$ to test.}
\KwOut{\textsc{Eliminated} or \textsc{Alive}, with a certificate set if eliminated.}
$W \gets w_x + r_x$ where $r_x = \sum_{j} g_{xj}$.\;
\lIf{$\exists\, i\ne x:\ w_i > W$}{\Return \textsc{Eliminated} (trivial, certificate $\{i\}$)}
Build the elimination network; $R \gets \sum_{i<j,\,i,j\ne x} g_{ij}$.\;
$(F,(S,T)) \gets \textsc{MaxFlowMinCut}(s,t)$.\;
\eIf{$\abs{F} = R$}{\Return \textsc{Alive}}{\Return \textsc{Eliminated}, certificate $A = \{i : i\in S\}$.\;}
\end{algorithm}

\begin{complexitybox}
For one team: $\Oh(n^2)$ game nodes, $\Oh(n^2)$ edges, max flow with Dinic in $\Oh(\abs{V}^2\abs{E}) = \Oh(n^6)$ worst case but vastly faster in practice (the graph is unit-ish and shallow). Testing all $n$ teams multiplies by $n$. Hoffman and Rivlin's classical condition sometimes lets you skip the flow, but the flow formulation is the clean general answer.
\end{complexitybox}

\begin{historybox}
The problem and a combinatorial criterion go back to Benjamin Schwartz (1966) and to Alan Hoffman and T.~J.~Rivlin (1967). The crisp max-flow formulation with the average-wins certificate became a textbook staple through Kevin Wayne's exposition. It is a favorite precisely because the reduction is so unobvious: ``can my team still win the pennant'' has no edges or capacities in its statement.
\end{historybox}

\section{Disjoint paths and Menger's theorem}
\label{ch18:menger}

The deepest classical use of flow is connectivity. How many \emph{edge-disjoint} paths run from $s$ to $t$? How many \emph{internally vertex-disjoint} ones? Menger's theorem answers both, and max-flow min-cut is the modern proof.

\begin{definition}[Disjoint paths and separators]
\label{ch18:def-menger}
Two $s$--$t$ paths are \emph{edge-disjoint} if they share no edge, \emph{internally vertex-disjoint} if they share no vertex except $s$ and $t$. An \emph{edge cut} is a set of edges whose removal leaves no $s$--$t$ path; a \emph{vertex separator} is a set of vertices (other than $s,t$) whose removal leaves no $s$--$t$ path.
\end{definition}

\begin{theorem}[Menger, edge form]
\label{ch18:thm-menger-edge}
In any graph, the maximum number of pairwise edge-disjoint $s$--$t$ paths equals the minimum number of edges in an $s$--$t$ edge cut.
\end{theorem}

\begin{theorem}[Menger, vertex form]
\label{ch18:thm-menger-vertex}
If $s,t$ are non-adjacent, the maximum number of internally vertex-disjoint $s$--$t$ paths equals the minimum size of an $s$--$t$ vertex separator.
\end{theorem}

\subsection{Edge form via unit capacities}

\begin{intuitionbox}
Give every edge capacity $1$. A flow of value $k$ is, by integrality of max flow, $k$ units that decompose into $k$ edge-disjoint paths (each unit reuses no edge, since capacity is $1$). A cut of capacity $k$ is literally a set of $k$ edges whose removal disconnects $s$ from $t$. Max-flow min-cut equates the two. That is the entire proof.
\end{intuitionbox}

\begin{proof}[Proof of Theorem~\ref{ch18:thm-menger-edge}]
Set $c(e)=1$ for all edges. By the integral flow theorem (max flow with integer capacities has an integral optimum), there is a max flow $f$ with $f(e)\in\{0,1\}$ and value $k=\abs{f}$. Decompose $f$ into paths and cycles; discard cycles (they do not change the value). Since each edge carries at most $1$ unit, no two of the resulting $s$--$t$ paths share an edge: $k$ edge-disjoint paths exist.

Conversely, any $k$ edge-disjoint paths give a feasible flow of value $k$ (one unit per path), so $\max f \ge$ (max edge-disjoint paths); the decomposition gives the reverse, hence equality between max flow and max edge-disjoint paths. Min-cut capacity, with unit capacities, counts edges, and a finite-capacity $s$--$t$ cut is exactly an $s$--$t$ edge cut. Max-flow min-cut closes the loop: max edge-disjoint paths $=\max f = \min$ cut $=$ min edge cut.
\end{proof}

\begin{center}
\begin{tikzpicture}[>=Stealth, every node/.style={font=\small,circle,draw,fill=nodeblue!15}, node distance=12mm]
  \node[fill=sourceblue!25] (s) {$s$};
  \node[above right=5mm and 16mm of s] (a) {$a$};
  \node[below right=5mm and 16mm of s] (b) {$b$};
  \node[right=16mm of a] (c) {$c$};
  \node[right=16mm of b] (d) {$d$};
  \node[fill=targetgreen!25, below right=5mm and 16mm of c] (t) {$t$};
  \draw[->] (s)--(a); \draw[->] (s)--(b);
  \draw[->] (a)--(c); \draw[->] (b)--(d);
  \draw[->] (a)--(d); \draw[->] (c)--(t); \draw[->] (d)--(t);
\end{tikzpicture}
\end{center}

\begin{labbox}
Trace the picture, all capacities $1$. Two edge-disjoint $s$--$t$ paths: $s\to a\to c\to t$ and $s\to b\to d\to t$. Can we get a third? Every $s$--$t$ path must use $s\to a$ or $s\to b$; there are only those two edges out of $s$, so at most $2$ edge-disjoint paths leave $s$. Indeed the cut $\{s\to a,\ s\to b\}$ has size $2$ and removing it isolates $s$. Max flow $=2=$ min cut. The diagonal $a\to d$ does not create a third disjoint path because the source still has out-degree $2$.
\end{labbox}

\subsection{Vertex form via vertex splitting}

For internally vertex-disjoint paths, paths may share edges only if they also share vertices, so the binding constraint is on vertices. Apply the vertex-split gadget of Section~\ref{ch18:vertex-cap} with every internal vertex capacity $1$.

\begin{focusbox}
\textbf{Vertex-disjoint network.} Split each vertex $v\notin\{s,t\}$ into $v_{\text{in}}\to v_{\text{out}}$ with capacity $1$; leave $s,t$ uncapacitated. Original edge $(u,v)$ becomes $(u_{\text{out}}, v_{\text{in}})$ with capacity $\infty$ (the vertex caps already do the limiting). Max flow $=$ max internally vertex-disjoint paths; a finite min cut must cut internal vertex edges and so equals a min vertex separator.
\end{focusbox}

\begin{proof}[Proof of Theorem~\ref{ch18:thm-menger-vertex}]
After splitting, every flow path traverses each internal vertex's unit-capacity internal edge $v_{\text{in}}\to v_{\text{out}}$, and the aggregate flow across that edge is at most $1$, so a flow of value $k$ decomposes into $k$ paths that pairwise share no internal vertex: $k$ internally vertex-disjoint paths. Conversely $k$ vertex-disjoint paths give a flow of value $k$. So max flow $=$ max vertex-disjoint paths.

A min cut in this network: edges crossing forward are either original $\infty$-edges (never in a finite cut) or internal $v_{\text{in}}\to v_{\text{out}}$ edges of capacity $1$. A finite min cut therefore selects a set of internal vertex-edges whose removal disconnects $s$ from $t$ -- that is exactly a vertex separator, of size equal to the cut capacity. Max-flow min-cut gives max vertex-disjoint paths $=$ min separator. Non-adjacency of $s,t$ ensures no direct $s\to t$ edge bypasses all internal vertices (which would make the separator undefined).
\end{proof}

\begin{labbox}
Take the diamond with edges $s\to a$, $s\to b$, $a\to t$, $b\to t$, and a cross edge $a\to b$. Internally vertex-disjoint $s$--$t$ paths: $s\to a\to t$ and $s\to b\to t$ are vertex-disjoint (share only $s,t$). A third would need a new internal vertex, but only $a,b$ exist, both used. The minimum separator is $\{a,b\}$, size $2$: deleting both kills every $s$--$t$ path. After splitting, the internal edges $a_{\text{in}}\to a_{\text{out}}$ and $b_{\text{in}}\to b_{\text{out}}$ (capacity $1$ each) form the min cut of capacity $2 = $ max flow. The cross edge $a\to b$ becomes $a_{\text{out}}\to b_{\text{in}}$ and cannot manufacture a third disjoint path, because any path through it still consumes both $a$'s and $b$'s unit capacity.
\end{labbox}

\begin{whybox}
Why does integrality matter so much here? The whole equivalence ``flow value $=$ number of disjoint paths'' rests on max flow having an \emph{integer} optimum when capacities are integers (proved in the max-flow chapter via augmenting paths preserving integrality). Without it, a fractional flow could split a unit across several edges and you could not read off disjoint paths. Unit capacities plus integrality is the engine behind every disjoint-paths and matching reduction.
\end{whybox}

\begin{connectionbox}
Menger is the local ($s$--$t$) version of global connectivity. The minimum over \emph{all} pairs $s,t$ of the min $s$--$t$ edge cut equals the global edge connectivity, computed in Chapter~11 (global min cut, Stoer--Wagner / Karger). The vertex form underlies $k$-connectivity testing and is the combinatorial heart of the max-flow proof of Hall's theorem and K\H{o}nig's theorem for bipartite matching (Chapter~15): a bipartite matching is just a unit-capacity $s$--$t$ flow, and K\H{o}nig's min-vertex-cover $=$ max-matching is Menger's vertex form on the bipartite flow network (Exercise~\ref{ch18:ex-konig}).
\end{connectionbox}

\begin{historybox}
Karl Menger proved his theorem in 1927, well before flows were formalized. Ford and Fulkerson's 1956 max-flow min-cut theorem subsumed it as the unit-capacity special case, and the modern view is that Menger, K\H{o}nig, Hall, Dilworth's chain-decomposition theorem, and max-flow min-cut are all the same duality wearing different clothes. Recognizing which costume a problem is wearing is the modeling skill this chapter is about.
\end{historybox}

\section{A reduction checklist}
\label{ch18:checklist}

\begin{focusbox}
When you suspect a problem is a flow/cut/matching, run this checklist.
\begin{enumerate}
\item \textbf{Binary choice with split penalties?} Min cut. Vertices $=$ items, source/sink $=$ the two labels, neighbor edges $=$ separation penalties (must be nonnegative: submodular).
\item \textbf{Prerequisites / ``take this forces that''?} Maximum-weight closure $\Rightarrow$ min cut, with $\infty$ on the forcing edges.
\item \textbf{Throughput through bottlenecks?} Max flow. Vertex limits $\Rightarrow$ split vertices. Minimum loads $\Rightarrow$ lower bounds.
\item \textbf{Assign items to bins under per-bin quotas?} Max flow: item$\to$bin edges, bin$\to t$ capacity $=$ quota (baseball is this).
\item \textbf{Disjointness / connectivity?} Unit capacities $\Rightarrow$ Menger. Vertex-disjoint $\Rightarrow$ split.
\end{enumerate}
\end{focusbox}

\begin{trapbox}
The single most common modeling bug is an \emph{orientation} or \emph{capacity-sign} error that makes the min cut measure the wrong thing -- e.g.\ putting profit on the sink side, or forgetting the $\infty$ on a forcing edge so the cut ``cheats'' by severing a prerequisite. Always verify the two-sided correspondence on a tiny instance by hand (as in every \textsf{Trace it yourself} box above) before trusting the solver. If your hand trace and the solver disagree, the model is wrong, not the solver (Exercise~\ref{ch18:ex-break}).
\end{trapbox}

\begin{aestheticbox}
Notice the same theorem -- max-flow min-cut -- powering five problems that share no surface vocabulary: mining schedules, image pixels, pennant races, network reliability, and graph connectivity. The reductions are not tricks; they are evidence that ``partition into two sides and pay for crossings'' and ``route as much as possible'' are among the deepest organizing ideas in combinatorial optimization. Learn to see the cut hiding inside a problem and a surprising fraction of optimization becomes a single subroutine call.
\end{aestheticbox}

\section{Exercises}
\label{ch18:exercises}

\begin{exercise}[Vertex-capacitated max flow]
\label{ch18:ex-vertexcap}
You are given a directed network where both edges and vertices have capacities. Describe the splitting construction precisely and prove that the resulting edge-capacitated max flow equals the original vertex-and-edge-capacitated max flow. What is the blow-up in $\abs{V}$ and $\abs{E}$?
\end{exercise}

\begin{exercise}[Lower bounds, end to end]
\label{ch18:ex-lower}
Run the lower-bound feasibility reduction of Algorithm~\ref{ch18:alg-lowerbound} on the following circulation: vertices $\{1,2,3\}$, edges $1\to2$ with $\ell=1,c=3$; $2\to3$ with $\ell=2,c=2$; $3\to1$ with $\ell=0,c=5$. Compute every $d(v)$, build the super-source/super-sink network, and decide feasibility by hand. Then exhibit a feasible flow or a certificate of infeasibility.
\end{exercise}

\begin{exercise}[Closure at the extremes]
\label{ch18:ex-closure-neg}
Suppose every vertex weight is negative. What is the maximum-weight closure, and what min cut does Theorem~\ref{ch18:thm-closure} produce? Now suppose every weight is positive: same questions. Use these extremes to sanity-check that an implementation handles the empty and full closures.
\end{exercise}

\begin{exercise}[Densest subgraph]
\label{ch18:ex-densest}
The density of a vertex subset $A$ in an undirected graph is $\abs{E(A)}/\abs{A}$ (edges inside $A$ over its size). Show that ``is there a subset of density $\ge g$?'' reduces to a maximum-weight closure / min-cut instance with a parameter $g$, and describe how a binary (or parametric) search over $g$ finds the densest subgraph. (Hint: make each edge a positive-weight ``project'' requiring both its endpoints, each endpoint a cost-$g$ prerequisite.)
\end{exercise}

\begin{exercise}[Segmentation hardness boundary]
\label{ch18:ex-seg-hard}
Prove that if some pairwise term \emph{rewards} disagreement -- i.e.\ a neighbor pair pays a negative $\lambda_{pq}$ when split -- then the energy is no longer the capacity of any $s$--$t$ cut with nonnegative capacities. Explain in one paragraph why this corresponds to non-submodularity and why min cut can no longer be used directly.
\end{exercise}

\begin{exercise}[Multiway labels]
\label{ch18:ex-multiway}
Three labels $\{R,G,B\}$, energy with data terms and a Potts smoothness ($\lambda$ paid per disagreeing neighbor pair). Show that the natural ``one source per label'' min-cut construction does \emph{not} in general give the optimum (find a 3-pixel counterexample), and explain how $\alpha$-expansion sidesteps this with a sequence of binary cuts.
\end{exercise}

\begin{exercise}[Baseball certificate]
\label{ch18:ex-baseball-cert}
Standings $w = (75, 71, 69, 63)$, remaining games $g_{12}=2, g_{13}=1, g_{14}=0, g_{23}=4, g_{24}=0, g_{34}=3$. Test team $4$ assuming it can reach $W=71$ in the best case. Build the elimination network for $x=4$, compute the max flow versus $R$, and -- if eliminated -- produce the subset-$A$ certificate from the min cut and verify $(w(A)+g(A))/\abs{A} > W$.
\end{exercise}

\begin{exercise}[Menger by hand]
\label{ch18:ex-menger}
For the graph with undirected unit-capacity edges $sa, sb, ab, at, bc, ct, ac$, find the maximum number of edge-disjoint $s$--$t$ paths and a minimum edge cut, and verify they are equal. Then redo it for internally vertex-disjoint paths and minimum vertex separators using the vertex-split network.
\end{exercise}

\begin{exercise}[K\H{o}nig from Menger]
\label{ch18:ex-konig}
Derive K\H{o}nig's theorem (in a bipartite graph, max matching $=$ min vertex cover) from Menger's vertex form (Theorem~\ref{ch18:thm-menger-vertex}) applied to the standard bipartite flow network ($s\to$ left $\to$ right $\to t$, unit capacities). Identify what an internally vertex-disjoint $s$--$t$ path is, and what a vertex separator is, in graph-matching terms.
\end{exercise}

\begin{exercise}[Project selection with mutual exclusion]
\label{ch18:ex-mutex}
Extend the project-selection model so that some pairs of projects are \emph{mutually exclusive} (at most one of the pair may be chosen). Show that a naive $\infty$-edge gadget cannot encode ``at most one,'' and discuss why mutual exclusion can make the problem NP-hard (relate it to the non-submodular / max-cut boundary). Give one special case where it stays polynomial.
\end{exercise}

\begin{exercise}[Break the model]
\label{ch18:ex-break}
Take the project-selection construction of Algorithm~\ref{ch18:alg-projsel} and \emph{remove} the $\infty$ on prerequisite edges, replacing it with a large finite capacity $M$. Construct an instance where the min cut now ``cheats'' by severing a prerequisite edge, returning an infeasible (non-closed) selection. State the smallest $M$ for which your instance still misbehaves, and explain the fix.
\end{exercise}
